\documentclass[conference]{IEEEtran}
\IEEEoverridecommandlockouts
% The preceding line is only needed to identify funding in the first footnote. If that is unneeded, please comment it out.
\usepackage{cite}
\usepackage{amsmath,amssymb,amsfonts}
\usepackage{algorithmic}
\usepackage{graphicx}
\usepackage{textcomp}
\usepackage{xcolor}
\def\BibTeX{{\rm B\kern-.05em{\sc i\kern-.025em b}\kern-.08em
    T\kern-.1667em\lower.7ex\hbox{E}\kern-.125emX}}

\newcommand{\pmtwofive}{\text{PM}_{2.5}}
\begin{document}

\title{Conference Paper Title}

\author{
  \IEEEauthorblockN{Mughilvannan Manivannan}
  \IEEEauthorblockA{\textit{Oslomet} \\ muman2187@oslomet.no} \and
  \IEEEauthorblockN{Ahamed Hasan}
  \IEEEauthorblockA{\textit{Oslomet} \\ mrdas1414@oslomet.no} \and
  \IEEEauthorblockN{Mrinmoy Das}
  \IEEEauthorblockA{\textit{Oslomet} \\ ahhas2658@oslomet.no}
}

\maketitle

\begin{abstract}
  Having a reliable and local AQI monitor is important in this day and age.
  Traditional methods to measure the AQI cannot be deployed in scale.
\end{abstract}

\section{Introduction}

\section{Brief Literature Review}

In the recent years there has been a explosion of research on AQI/$\pmtwofive$
detection from images. They promise a cheaper alternative to government grade
air quality monitoring stations which are sparse and expensive to maintain using
smartphones, traffic cameras and drones. In this section we will see how the
technology has evolved in the past 3 years.

\subsection{Phase 1: Physics of light}

The filed started with atmospherics physics and how light interact with them.
The foundation of how the light gets absorbed and scattered are governed by
Beer-Lambert law and Koschmieder's visibility model. Airborne aerosols absorb
and scatter ambient light reducing contrast, washing out colour saturation and
blurring distant objects. Early research built systems around handcrafted
optical descriptors and Image Quality Assessment metrics.

As systematically reviewed by Marzagui et al.
(2024)\cite{marzaguiReviewImageQuality2024}, early algorithms tracked
degradation using the Dark Channel Prior (DCP), transmission matrices,
structural similarity, and natural scene statistics. Yadav et al.
(2025)\cite{yadavRoadsideTreeImageBased2025} explored roadside tree imagery
across Gurugram, India, discovering that while standard texture descriptors like
HOG and SIFT failed, calculating the physical atmospheric extinction coefficient
yielded a 68\% classification accuracy. Chakraborty et al.
(2024)\cite{chakrabortyCanImageTell2024} pushed optical modelling further: by
applying a Pyramid Fusion dehazing algorithm to create a synthetic reference
image, they computed the Structural Similarity Index Measure (SSIM) across
distant Regions of Interest (RoIs) at known optical depths, reaching an $R^2$ of
0.93 with an Extra Trees Regressor.

However, handcrafted features remained brittle: they required fixed camera
angles, struggled with cloudy weather, and failed to differentiate natural
humidity from industrial smog.

\subsection{Phase 2: Deep Learning and Mobile Efficiency}

The next shift was seen when the Deep Learning model got more effective in
identifying the visual clues from the raw pixels using CNN. This replaced the
need for hand crafted features.

\begin{itemize}
  \item \textbf{Model Benchmarking:} Kaushal et al.
        (2024)\cite{kaushalEvaluationDeepLearning2024} demonstrated the
        advantage of residual learning by showing that ResNet-50 reached 97.87\%
        accuracy on Beijing outdoor imagery, outperforming VGG-19. Verma et al.
        (2025)\cite{vermaImagebasedAirPollution2025} evaluated a separable
        Conv2D CNN against an HOG-LSTM architecture on 12,240 images across
        India and Nepal, proving the CNN significantly outperformed temporal
        feature extractors (96.49\% vs. 85.87\% accuracy).
  \item \textbf{Lightweight Architectures:} Moving algorithms to mobile and
        Internet of Things (IoT) hardware requires lean neural footprints.
        Kumaar.B and Deepak (2025) found that the lightweight MobileNetV2
        (83.65\% accuracy) generalized better and ran faster than the deeper
        ResNet152V2 (78.79\%). Wang et al. (2024) introduced MDNet, using
        multiscale dilated convolutions to expand the network's receptive field
        to capture sky haze while cutting parameters by 86.7\% (achieving 94.2\%
        accuracy with only 4.8M parameters). Similarly, Zhang et al. (2024)
        enhanced ResNet18 with Inverted Bottleneck blocks and ParNet channel
        attention, cutting parameter size by 21.68\% while boosting accuracy to
        90.65\%.
  \item \textbf{Generative Communication:}
\end{itemize}

\section{Extended Literature Review}

\bibliographystyle{IEEEtran}
\bibliography{references}
\end{document}
